program 08 \documentclass[12pt]{article}

% Package for mathematical symbols and theorem environments
\usepackage{amsmath, amsthm, amssymb}

% Define theorem-like environments
\newtheorem{theorem}{Theorem}[section]
\newtheorem{lemma}[theorem]{Lemma}
\newtheorem{corollary}[theorem]{Corollary}

% Definition style
\theoremstyle{definition}
\newtheorem{definition}[theorem]{Definition}

\begin{document}

\title{Presentation of Numbered Theorems, Definitions, Corollaries, and Lemmas}
\author{}
\date{}
\maketitle

\section{Mathematical Statements}

\begin{definition}
A prime number is a natural number greater than 1 that has no positive divisors other than 1 and itself.
\end{definition}

\begin{theorem}
Every prime number greater than 2 is an odd number.
\end{theorem}

\begin{proof}
Assume that a prime number greater than 2 is even. Then it must be divisible by 2.

Since it is greater than 2, it has at least three divisors: 1, 2, and itself.

This contradicts the definition of a prime number, which must have exactly two positive divisors.

Therefore, every prime number greater than 2 is odd.
\end{proof}

\begin{lemma}
If an integer is divisible by 4, then it is also divisible by 2.
\end{lemma}

\begin{proof}
Let the integer be \( n = 4k \), where \( k \) is an integer.

Then,
\[
n = 4k = 2(2k)
\]

Since \( 2k \) is also an integer, \( n \) is divisible by 2.
\end{proof}

\begin{corollary}
The number 16 is divisible by 2.
\end{corollary}

\begin{proof}
Since \(16 = 4 \times 4\), it is divisible by 4.

By the lemma, every number divisible by 4 is also divisible by 2.

Hence, 16 is divisible by 2.
\end{proof}

\end{document}